%

%\documentclass[11pt,a4paper,twoside]{article}

\usepackage[utf8]{inputenc}

\usepackage[T1,T2A]{fontenc}

\usepackage[english.ancient,russian.ancient]{babel}

\usepackage{graphicx}

\usepackage{array}

\usepackage[doi]{samgtu-bib}

\usepackage{samgtu}

\usepackage{samgtu-name}

\usepackage{mathtools}

\mathtoolsset{showonlyrefs}

\usepackage{desclist}

\usepackage{euscript}

\usepackage{mathrsfs} 

\usepackage{multicol}

\usepackage{mathrsfs}

\usepackage{cite}

\usepackage{cmap}

\usepackage{qrcode}

\allowdisplaybreaks[4]

\usepackage[nodayofweek]{datetime}

\usepackage[thinlines]{easybmat}



\renewcommand{\JournalYear}{2018}

\renewcommand{\JournalVolume}{22}

\renewcommand{\JournalNumber}{4}



\newdate{JournalDateSign}{29}{12}{2018}% Дата подписания Вестника в печать

\renewcommand{\JournalDateSign}{\displaydate{JournalDateSign}}

\newdate{JournalDatePrint}{16}{01}{2019}% Дата выхода Вестника из печати

\renewcommand{\JournalDatePrint}{\displaydate{JournalDatePrint}}

%\renewcommand{\JournalUPL}{16.56} % Усл. печ. л.

%\renewcommand{\JournalUIL}{16.53} % Уч.-изд. л.

\renewcommand{\JournalUPL}{15.24} % Усл. печ. л.

\renewcommand{\JournalUIL}{15.21} % Уч.-изд. л.

\renewcommand{\JournalRegNumber}{220/18} % Регистрационный номер

\renewcommand{\JournalOrderNumber}{2} % Номер заказа



%\renewcommand{\JournalRMonth}{Март} % Месяц выхода Вестника

%\renewcommand{\JournalEMonth}{March} % Месяц выхода Вестника

%\renewcommand{\JournalRMonth}{Июнь} % Месяц выхода Вестника

%\renewcommand{\JournalEMonth}{June} % Месяц выхода Вестника

%\renewcommand{\JournalRMonth}{Сентябрь} % Месяц выхода Вестника

%\renewcommand{\JournalEMonth}{September} % Месяц выхода Вестника

\renewcommand{\JournalRMonth}{Декабрь} % Месяц выхода Вестника

\renewcommand{\JournalEMonth}{December} % Месяц выхода Вестника



\newcommand{\totop}[1]{\raisebox{6pt}[1pt][2pt]{#1}}

\newcommand{\tottop}[1]{\raisebox{12pt}[1pt][2pt]{#1}} 

\newcommand{\totttop}[1]{\raisebox{18pt}[1pt][2pt]{#1}} 

\newcommand{\tottttop}[1]{\raisebox{24pt}[1pt][2pt]{#1}} 

\newcommand{\totttttop}[1]{\raisebox{30pt}[1pt][2pt]{#1}} 

\newcommand{\tobot}[1]{\raisebox{-6pt}[1pt][2pt]{#1}} 

\newcommand{\tobbot}[1]{\raisebox{-12pt}[1pt][2pt]{#1}} 

\newcommand{\tobbbot}[1]{\raisebox{-18pt}[1pt][2pt]{#1}}



\graphicspath{{./00PIC/}}

%

%

%

%

%\pdfcrop % обрезка pdf для размещения на math-net.ru

%\draft % На корректуре

%

%

%

\FirstPage{388}

\LastPage{397}

%

%

%\begin{document}



\newpage





\selectlanguage{russian}



\phantomsection



\addcontentsline{toc}{section}{Редакционная политика журнала}

\addcontentsline{etoc}{section}{Editorial Policy}







\vsgtu{}

\chead[\fancyplain{}{\selectlanguage{russian}\therutitle}]{\fancyplain{}{\selectlanguage{russian} \therutitle}}

\rhead[]{\fancyplain{\RuJournal}{}}

\lhead[\fancyplain{\RuJournal}{}]{}

\cfoot[]{}

\lfoot[\thepage]{}

\rfoot[]{\thepage}

\renewcommand{\plainheadrulewidth}{0.7pt}

\renewcommand{\headrulewidth}{0.7pt}

\pagestyle{fancyplain}

\thispagestyle{plain}

\normalfont



$\;$



\bigskip



\rutitle{Редакционная политика журнала}





\begin{center}



{\large \bf Редакционная политика журнала «Вестник Самарского государственного технического университета. Серия:~Физико-математические науки»\footnote{Настоящая редакционная политика опубликована на сайте журнала и доступна по ссылке \url{http://www.mathnet.ru/rus/vsgtu/edpolicy}. Всегда проверяйте наличие новой редакции политики журнала на сайте!}}



\smallskip



{\small Редакционная коллегия журнала}



\smallskip



{\small (01 июня 2017 г.)}



\end{center}





\renewcommand{\abstractname}{Общие сведения о журнале}





\small



\begin{abstract}

Журнал «Вестник Самарского государственного технического университета. Серия: Физико"=математические науки» (параллельное название на английском языке "--- Journal of Samara State Technical University, Ser. Physical and Mathematical Sciences) является рецензируемым научным журналом Самарского государственного технического университета и~издаётся с~1996~г.



Самарский государственный технический университет сегодня "--- это один из одиннадцати ведущих опорных университетов России, который является крупнейшим научно-образовательным центром Поволжья с~более чем 100-летней историей.



Журнал долгое время позиционировал себя как издание, предназначенное для публикации новых научных знаний, полученных в российских научных школах. Однако в~настоящее время журнал ставит своей целью открытое распространение научных знаний среди российских и зарубежных учёных, поэтому он также ориентируется и~на зарубежных учёных, работающих в приоритетных научных направлениях Самарского государственного технического университета.



Журнал зарегистрирован в~Федеральной службе по надзору в сфере связи, информационных технологий и массовых коммуникаций (Роскомнадзор), свидетельство \href{https://rkn.gov.ru/mass-communications/reestr/media/?id=83010}{ПИ № ФС 77–66685 от 27.07.2016}. С~2011 года журнал выходит ежеквартально; объём номера "--- 200~c.; язык публикаций "--- русский, английский. Журнал издаётся в~печатной и~электронной формах.



Редакция журнала принимает и оценивает рукописи представленных статей независимо от расы, пола, национальности, происхождения, гражданства (подданства), рода занятий, места работы и проживания автора, а также от его политических, философских, религиозных и иных взглядов.



Представляемая в~журнал рукопись статьи должна быть законченным научным исследованием, нигде ранее не публиковавшимся и~не представленным к~публикации в~других изданиях.



Рукопись статьи должна содержать новые научные результаты по приоритетным направлениям Самарского государственного технического университета, таким как «Дифференциальные уравнения и математическая физика», «Механика деформируемого твёрдого тела», «Математическое моделирование, численные методы и~комплексы программ».



Журнал издается на средства издателя. Все публикации в~журнале бесплатны. Все публикации в~электронном виде распространяются бесплатно.



Целевую аудиторию журнала составляют учёные и исследователи, чьи научные интересы лежат в указанных направлениях: «Дифференциальные уравнения и математическая физика», «Механика деформируемого твёрдого тела», «Математическое моделирование, численные методы и комплексы программ».



\end{abstract}





\bigskip





\noindent

\textbf{1. Цели и задачи}



\begin{enumerate}



\item[1.] Журнал долгое время позиционировал себя как издание, предназначенное для публикации новых научных знаний, полученных в российских научных школах. Однако в~настоящее время журнал ставит своей целью открытое распространение научных знаний среди российских и зарубежных учёных, поэтому он также ориентируется и на зарубежных учёных, работающих в приоритетных научных направлениях Самарского государственного технического университета.



\item[2.] Журнал публикует оригинальные статьи и заказные обзоры по следующим направлениям: «Дифференциальные уравнения и~математическая физика», «Механика деформируемого твёрдого тела», «Математическое моделирование, численные методы и комплексы программ».



\item[3.] Журнал особое внимание уделяет следующим задачам:

\begin{itemize}

\item привлечение высококвалифицированных авторов, редакторов и рецензентов по профилю журнала;

\item обеспечение высокого научного уровня рецензирования всех поступающих материалов;

\item соответствие политики журнала международным стандартам  (стандартам \href{http://publicationethics.org/}{COPE}).

\end{itemize}



\item[4.] Целевую аудиторию журнала составляют учёные и~исследователи, чьи научные интересы лежат в~указанных направлениях: «Дифференциальные уравнения и математическая физика», «Механика деформируемого твёрдого тела», «Математическое моделирование, численные методы и комплексы программ». 



\end{enumerate}



\noindent

\textbf{2. Политика разделов журнала}





\begin{enumerate}



\item[1.] Журнал публикует оригинальные статьи, заказные обзоры и~краткие сообщения по следующим направлениям: «Дифференциальные уравнения и математическая физика», «Механика деформируемого твёрдого тела», «Математическое моделирование, численные методы и комплексы программ».

Журнал может публиковать и другие материалы (биографии, письма к редактору, комментарии и т.п.).



\item[2.] 

Журнал не публикует материалы рекламного характера.

\item[3.] При наличии в редакционном портфеле нескольких работ одного автора в~текущем номере может быть опубликована только одна работа (по выбору редакционной коллегии), вторая "--- в~следующем номере и~так далее. В~исключительных случаях редакционная коллегия может принять решение об опубликовании не более двух работ одного автора. 

\item[4.] 

Весь контент журнала индексируется в~базах данных.



\end{enumerate}





\textbf{2.1. Научные статьи}



\hypertarget{sectionPoliciesOriginalArticles}{}



\begin{enumerate}



\item[1.] В научных статьях должны быть представлены оригинальные исследования с~чётко сформулированными гипотезами, целями и~задачами. Статьи должны содержать новые подходы и новое понимание рассматриваемых проблем.



\item[2.] Все рукописи должны быть подготовлены с помощью \LaTeXe\ (при этом должен быть использован \href{https://www.overleaf.com/read/bzytwtpwfzjj}{наш шаблон}) и отправлены редактору с помощью \href{http://www.mathnet.ru/php/esubmission.phtml?jrnid=vsgtu&wshow=snote&option_lang=rus}{онлайн-службы представления рукописей}.



\item[3.]

Как правило, текст статьи не должен превышать 7500 слов или 30 страниц (без библиографического списка). Статья может включать 8 рисунков или таблиц. Статья должна содержать не менее 10 и не более 50 ссылок. Некоторые рукописи могут содержат больше слов и/или ссылок (по разрешению редактора), другие – меньше.



\item[4.] Все научные статьи должны быть структурированы (например, в формате IMRaD: введение, методы, результаты и дискуссия).



\item[5.] При подготовке рукописи авторы должны пользоваться \href{http://www.mathnet.ru/rus/vsgtu/authornotes#manuscriptPreparingRequirements}{руководством для написания статей}.



\item[6.] Представляемая в~журнал рукопись статьи должна быть законченным научным исследованием, нигде ранее не публиковавшимся и~не представленным к~публикации в~других изданиях. 



\item[7.] Все научные статьи отправляются на рецензирование.



\end{enumerate}





{\bf 2.2. Обзоры}



\begin{enumerate}



\item[1.] Обзор "--- попытка одного или нескольких авторов обобщить текущее состояние исследования по определенной теме (предметной области). В~идеале автор ищет все, что относится к~теме обзора, а~затем сортирует все это и формирует представление о~текущем «состоянии дел». 



\item[2.] Обзор должен содержать информацию об:

\begin{itemize}

\item основных научных публикациях по теме обзора;

\item последние крупные достижения и открытия;

\item значительные пробелы в исследованиях;

\item основные спорные вопросы по теме обзора;

\item основные научные вопросы и перспективы их решения.

\end{itemize}



\item[3.]  Обзоры могут сильно различаться по длине. Обзор может варьироваться от 8000 до 40000 слов или от~30 до 90 страниц (включая библиографический список, таблицы, рисунки).



\item[4.]  Для публикации принимаются обзоры только  от экспертов в~конкретной предметной области.



\item[5.]  Все обзоры отправляются на рецензирование.



\end{enumerate}





{\bf 2.3. Краткие сообщения}



\begin{enumerate} 

\item[1.] Краткое сообщение предназначено для печати небольшого, но независимого результата, вносящего значительный вклад в~область исследования.

\item[2.] Краткое сообщение не должно содержать предварительных результатов, если только эти результаты не представляют исключительный интерес и особенно актуальны.

\item[3.] К подаче и~структуре краткого сообщения предъявляются такие же требования, как и~к~\hyperlink{sectionPoliciesOriginalArticles}{научным статьям}.

\item[4.] Размер краткого сообщения не должен превышать 2500 слов или 10 страниц (без библиографического списка). Оно может включать два рисунка или две таблицы. Краткое сообщение должно содержать не менее 8 ссылок.



\item[5.] Краткие сообщения также отправляются на рецензирование.



\end{enumerate}



\pagebreak



{\bf 2.4. Другие материалы (биографии, письма к редактору, комментарии}



{\bf и т.п.)}



\begin{enumerate}

\item[1.] Размер таких материалов не должен превышать 2500 слов или 10 страниц.

\item[2.] Такие материалы не рецензируются. 

\end{enumerate}



\smallskip



\noindent

{\bf 3. Периодичность}



\smallskip



С~2011 года журнал выходит ежеквартально; объём номера "--- 200~c.; язык публикаций "--- русский, английский.



\smallskip



\smallskip



\smallskip



\noindent

{\bf 4. Политика свободного доступа}



\begin{enumerate}

\item[1.] Журнал сразу предоставляет открытый доступ к~своему контенту исходя из следующего принципа: свободный открытый доступ к результатам исследований способствует увеличению глобального обмена знаниями.



\item[2.] Все публикации журнала в электронном виде распространяются бесплатно и~без ограничений.



\item[3.] Весь контент журнала распространяется на условиях лицензии \href{https://creativecommons.org/licenses/by/4.0/deed.ru}{Creative Co\-m\-mons Attribution 4.0 Inter\-na\-tio\-nal} (\href{https://creativecommons.org/licenses/by/4.0/deed.ru}{CC BY 4.0}; \url{https://creativecommons.org/licenses/by/4.0/deed.ru}).



\item[4.] По решению редакционной коллегии и издателя весь контент журнала, принятый к публикации и/или опубликованный до 01 января 2017, также распространяется по лицензии \href{https://creativecommons.org/licenses/by/4.0/deed.ru}{CC BY 4.0}. 



\item[5.] Если авторы опубликованного ранее контента не согласны с условиями лицензии \href{https://creativecommons.org/licenses/by/4.0/deed.ru}{CC BY 4.0}, то просим их незамедлительно сообщить об этом в редакционную коллегию по e-mail \href{mailto:vsgtu@samgtu.ru}{vsgtu@samgtu.ru}.



\end{enumerate}



\noindent

{\bf 5. Авторские права и политика архивирования}



\smallskip



Авторы, публикующиеся в данном журнале, соглашаются со следующим: 

\begin{enumerate}

\item[1.]  Весь контент журнала распространяется по лицензии \href{https://creativecommons.org/licenses/by/4.0/deed.ru}{CC BY 4.0}. Согласно \href{https://creativecommons.org/licenses/by/4.0/deed.ru}{CC BY 4.0} авторы сохраняют права на свои статьи, но при этом разрешают всем беспрепятственно скачивать, повторно использовать, перепечатывать, изменять, распространять и/или копировать их при условии упоминания их авторства. На всё вышеперечисленное разрешения от авторов или издателя не требуется.



\item[2.]  Авторы сохраняют право заключать отдельные контрактные договорённости, касающиеся неэксклюзивного распространения версии работы в опубликованном здесь виде (например, размещение ее в институтском хранилище, публикацию в книге, перевод на другой язык), со ссылкой на её оригинальную публикацию в~этом журнале.



\item[3.] Согласно классификации \href{http://www.sherpa.ac.uk/romeo/index.php}{SHERPA/RoMEO} журнал относится к так называемым «зеленым» журналам. Авторам разрешено архивировать препринты и постпринты своих работ. Авторы имеют право размещать свою работу в сети Интернет (например, в~институтском хранилище или персональном сайте) до и~во время процесса рассмотрения её в~редакционной коллегии журнала.



\end{enumerate}



\smallskip



\noindent

{\bf 6. Рецензирование}



\begin{enumerate}



\item[1.] Сначала рукопись читается и~проверяется редактором на соответствие руководству для авторов (первоначальная оценка рукописи).



\item[2.]  Несоблюдение изложенных в \href{http://www.mathnet.ru/rus/vsgtu/authornotes#manuscriptPreparingRequirements}{правилах для авторов} критериев может привести к возврату рукописи для исправления без рецензирования.



\item[3.]  Статьи, которые считаются подходящими, направляются как минимум двум независимым рецензентам для оценки научного качества работы.



\item[4.]  Все рецензенты являются признанными специалистами по тематике рецензируемых материалов и~имеют в~течение последних трёх лет публикации по тематике рецензируемой статьи.



\item[5.]  В~журнале действует процесс одностороннего слепого рецензирования или двойного слепого рецензирования, если это возможно сделать для конкретной статьи.



\item[6.] При подготовке рецензии рецензенты обязаны пользоваться инструкциями журнала по рецензированию и~этическими нормами для рецензентов. 



\item[7.] Рецензии хранятся в~редакции журнала в~течение пяти лет.



\item[8.] Рукописи \href{https://goo.gl/wDhKbS}{оцениваются} в соответствии со следующими аспектами:

\begin{itemize}

\item Представлены ли в документе новые концепции (теории), идеи или данные? 

\item Являются ли научные методы и допущения обоснованными и четко обозначенными? 

\item Являются ли результаты достаточными, чтобы обосновать выводы? 

\item Является ли описание экспериментов и~расчётов достаточно полным и точным, чтобы они могли быть воспроизведены другими учеными (возможна ли прослеживаемость результатов)?

\item Придают ли авторы должное внимание аналогичным работам, чётко ли указывается их собственный новый/оригинальный вклад в~развитие проблемы?

\item Насколько точно отражает название статьи её содержание?

\item Отражает ли аннотация основные идеи и результаты исследования?

\item Является ли рукопись статьи хорошо структурированной и~понятной?

\item Насколько хорош и точен язык статьи?

\item Правильно ли определены и используются математические формулы, символы, аббревиатуры и единицы?

\item Должны ли быть какие-либо части документа (текст, формулы, рисунки, таблицы) разъяснены, сокращены, объединены или устранены?

\item Удовлетворителен ли библиографический список?

\end{itemize}



\item[9.] Все вышеперечисленные вопросы будут оцениваться рецензентами в соответствии со следующими оценками: выдающиеся, хорошие, справедливые, посредственные или плохие.



\item[10.] Статья может быть: принята (без дальнейшей доработки); направлена на доработку; отклонена; отклонена (по причине очень предварительных результатов); отклонена (не соответствует тематике журнала).



\item[11.] Если статья направлена на доработку, она должна быть снова отправлена на рассмотрение тем же рецензентам.



\item[12.] Решение о~публикации принимается редакционной коллегией журнала на основании экспертных оценок рецензентов с~учётом соответствия представленных материалов тематической направленности журнала, их научной значимости и~актуальности.

\item[13.] Авторы рукописи имеют право рекомендовать рецензентов "--- специалистов в~данной области (не более двух персон).

\item[14.]  Выбор рецензентов редакционная коллегия оставляет за собой.

\item[15.] Редакционная коллегия не привлекают к~работе рецензентов, которые могут иметь конфликт интересов. Авторы могут сообщать имена людей, назначение которых рецензентами нежелательно из-за возможности возникновения конфликта интересов, объяснив свою позицию.



\end{enumerate}



\smallskip



\noindent

{\bf 7. Индексирование}



\begin{enumerate}

\item[1.] Журнал включен в  \href{http://elibrary.ru/rsci_about.asp}{Russian Science Citation Index} на платформе \href{http://wokinfo.com/products_tools/multidisciplinary/rsci/}{Web of Sci\-ence}. 



\item[2.] Журнал индексируется в реферативных базах данных \href{http://www.viniti.ru}{ВИНИТИ РАН}.

\item[3.] Статьи журнала индексируются также в \href{http://scholar.google.com/scholar?q=Vestn.+Samar.+Gos.+Tekhn.+Univ.+Ser.+Fiz.-Mat.+Nauki}{Scholar.Google.com}, \href{https://goo.gl/zrvH0K}{OCLC WorldCat}, Bielefeld Academic Search Engine (BASE), Open Access Infrastructure for Re\-search in Europe (OpenAIRE), Research Papers in Economics (RePEc), Socionet, \href{http://cyberleninka.ru/journal/n/vestnik-samarskogo-gosudarstvennogo-tehnicheskogo-universiteta-seriya-fiziko-matematicheskie-nauki}{СyberLeninka.ru}, \href{http://goo.gl/EeAvfJ}{Math-Net.ru}, \href{https://elibrary.ru/title_about.asp?id=5784}{eLibrary.ru}. 

\item[4.] Журнал интегрирован в поисковые системы 

\href{http://search.crossref.org/}{CrossRef} и \href{http://search.crossref.org/funding}{FundRef}.

\item[5.] Информация о журнале опубликована в ULRICH’S Periodical Directory. 



\end{enumerate}



\smallskip

\smallskip



\noindent

{\bf 8. Этика научных публикаций}\footnote{Настоящие положения в полном объёме взяты с сайта журнала «Вестник ПНИПУ. Механика» (\url{http://vestnik.pstu.ru/mechanics/toauthors/ethics/})}



\smallskip



\smallskip



\smallskip



\noindent

{\bf 8.1. Общие положения}



\begin{enumerate}

\item[\bf 1.1.] Настоящей политикой устанавливаются стандарты этического поведения для сторон, участвующих в процессе публикации: авторов, редакции, рецензента, издателя. Нижеперечисленные стандарты основаны на общепринятой и существующей политике журнала и издателя.

\item[\bf 1.2.]  Федеральное государственное бюджетное образовательное учреждение высшего образования «Самарский государственный технический университет», как учредитель и издатель научного журнала «Вестник Самарского государственного технического университета. Серия: Физико-математические науки», принимает на себя обязательства по контролю за всеми этапами публикации статей и признает свои этические и другие обязательства, связанные с опубликованием статей.

\end{enumerate}





\smallskip



\noindent

{\bf 8.2. Этические стандарты, предъявляемые к авторам публикаций}





\begin{enumerate}

\item[\bf 2.1.] {\bf Стандарт доступа к исходным данным исследования и их хранения.}

Автор обязан представить исходные материалы (данные) исследования по требованию редакции, если в статье не приводятся исходные данные, и должен быть готов предоставить публичный доступ к ним. Автор должен хранить эти данные в течение разумного времени после публикации для возможности их воспроизведения и проверки.



\item[\bf 2.2.] {\bf Стандарт оригинальности (недопустимости плагиата и самоплагиата).}

Автор представляет в редакцию для рассмотрения статью, содержащую результаты оригинального исследования. Если автор в статье использовал работы или включает в свою статью фрагменты из работ (цитаты) других лиц, то такое использование должно быть надлежащим образом оформлено путем указания оригинального источника в библиографическом списке к статье. Плагиат, равно как и автоплагиат, в любой форме является неэтичным и неприемлемым поведением автора.



\item[\bf 2.3.] {\bf Стандарт однократности публикации.}

Автор представляет в редакцию рукопись статьи, ранее не публиковавшейся и не переданной в редакции других журналов. Подача рукописи одновременно в несколько журналов является неэтичной и неприемлема. % Это же касается перевода статьи на иностранный язык.



\item[\bf 2.4.] {\bf  Стандарт подтверждения источников.}

Автор обязуется правильно указать научные и иные источники, которые он использовал в ходе исследования и которые оказали существенное влияние на результаты исследования, в библиографическом списке. Информация, полученная из неофициальных (частных) источников, не должна использоваться при оформлении научной статьи.



\item[\bf 2.5.] {\bf  Стандарт авторства рукописи статьи.}

Все лица, внесшие значительный вклад в получение результатов исследования, должны быть указаны в качестве соавторов статьи. Список авторов должен быть ограничен лишь этими лицами.

Автор, представляющий редакции рукопись, гарантирует, что им указаны все соавторы, что все они видели и одобрили окончательный вариант рукописи и согласны с её представлением в редакцию научного журнала «Вестник Самарского государственного технического университета. Серия: Физико-математические науки» для публикации.



\item[\bf 2.6.] {\bf  Стандарт раскрытия конфликта интересов со стороны автора.}

Авторы должны раскрывать конфликты интересов, которые могут повлиять на оценку и интерпретацию их рукописи. Все источники финансовой поддержки проекта (гранты, госпрограммы, проекты и проч.) должны быть раскрыты и в обязательном порядке указаны в рукописи.



\item[\bf 2.7.] {\bf  Стандарт исправления ошибок в опубликованных работах.}

Если автор обнаружит существенную ошибку или неточность в уже опубликованной статье, то он обязан незамедлительно уведомить об этом редакцию и содействовать ей в исправлении ошибки. Если редакция узнает об ошибке от третьих лиц, то автор обязан незамедлительно устранить ошибку или представить доказательства ее отсутствия.



\end{enumerate}



\smallskip



\noindent

{\bf 8.3. Этические стандарты, предъявляемые к редакции}





\begin{enumerate}

\item[\bf 3.1.] {\bf  Стандарт принятия решения об опубликовании статьи.}

Редакция журнала принимает решение о том, какие из статей, представленных в редакцию, должны быть опубликованы, на основании результатов проверки на предмет выполнения требований к оформлению и результатов рецензирования.

При принятии решения о публикации рукописи редакция руководствуется политикой журнала и не допускает публикацию статей с признаками клеветы, оскорбления, плагиата или нарушения авторских прав. Окончательное решение о публикации статьи или об отказе в таковой принимается главным редактором журнала.



\item[\bf 3.2.] {\bf  Стандарт равенства всех авторов.}

Редакция оценивает рукописи представленных статей независимо от расы, пола, национальности, происхождения, гражданства (подданства), рода занятий, места работы и проживания автора, а также от его политических, философских, религиозных и иных взглядов.

\item[\bf 3.3.] {\bf   Стандарт конфиденциальности.}

Редакция обязуется не раскрывать информацию о представленной рукописи никому, кроме автора, рецензента (при одностороннем слепом рецензировании), а при необходимости "--- издателя.



\item[\bf 3.4.] {\bf  Стандарт раскрытия конфликта интересов со стороны редакции.}\linebreak

Редакция гарантирует, что материалы рукописи, отклоненной от публикации, не будут использоваться в собственных исследованиях членов редколлегии без письменного согласия автора. Редколлегия будет отказываться от рассмотрения рукописи, если она состоит в каких-либо конкурентных отношениях с автором или организацией, связанной с результатами исследования, либо имеется иной конфликт интересов. Редакция обязуется требовать от всех участников процесса опубликования статьи раскрытия конкурирующих интересов.



\item[\bf 3.5.] {\bf  Стандарт рассмотрения претензий на неэтичное поведение автора.}

Редакция оперативно рассматривает каждую претензию на неэтичное поведение авторов рукописей и уже опубликованных статей независимо от времени её получения. Редакция обязуется принимать адекватные разумные меры в отношении таких претензий. Если доводы претензии подтвердятся, то редакция вправе отказаться от публикации статьи, прекратить дальнейшее сотрудничество с автором, опубликовать соответствующее опровержение, а также принять иные необходимые меры для дальнейшего пресечения неэтичного поведения данного автора.





\end{enumerate}



\smallskip



\noindent

{\bf 8.4. Этические стандарты, предъявляемые к рецензентам}



\begin{enumerate}

\item[\bf 4.1.] {\bf Стандарт вклада рецензента в редакционные решения.}

Предоставленная рецензентом экспертная оценка рукописи способствует принятию редакционных решений, а также помогает автору улучшить рукопись. Решение о принятии рукописи к публикации, возвращении ее автору на доработку или отклонении от публикации принимается редколлегией на основании результатов рецензирования.



\item[\bf 4.2.] {\bf Стандарт сроков рецензирования.}

Рецензент обязан предоставить рецензию в сроки, указанные редакцией. Если рассмотрение рукописи и подготовка рецензии в эти сроки невозможны, то рецензент должен уведомить об этом редакцию.



\item[\bf 4.3.] {\bf Стандарт конфиденциальности со стороны рецензента.}

Рукопись статьи, представленная на рецензирование, должна рассматриваться как конфиденциальный документ, независимо от избранной журналом формы рецензирования. Рецензент вправе показывать её или обсуждать с другими лицами только с разрешения главного редактора.

\item[\bf 4.4.] {\bf  Стандарт объективности рецензии.}

Рецензент обязуется проводить экспертную оценку рукописи объективно. Личная критика автора рецензентом недопустима. Все выводы рецензента должны быть строго аргументированы и снабжены ссылками на авторитетные источники.



\item[\bf 4.5.] {\bf  Стандарт подтверждения источников.}

Рецензенты должны указать при наличии работы, которые оказали влияние на результаты исследования, но не были приведены автором. Рецензент обязан обратить внимание редакции на существенное сходство или совпадение между рассматриваемой рукописью и ранее опубликованной другой работой, о которой известно рецензенту.



\item[\bf 4.6.] {\bf  Стандарт раскрытия конфликта интересов.}

Рецензент не может использовать материалы неопубликованной рукописи в своих собственных исследованиях без письменного согласия автора. Рецензент обязан отказаться от рассмотрения рукописи, в связи с которой у него возникает конфликт интересов из-за конкурентных, совместных или иных отношений с автором либо организацией, имеющей отношение к рукописи. 



\end{enumerate}



\smallskip



\pagebreak



\noindent

{\bf 9. Стоимость публикации}



\begin{enumerate}

\item[1.]

Журнал издается на средства издателя.

\item[2.]

Все публикации в журнале бесплатны (нет сборов за публикацию и подготовку статьи).

\end{enumerate}





\noindent

{\bf 10. Политика раскрытия и конфликты интересов}



\begin{enumerate}

\item[1.]

Неопубликованные данные, полученные из представленных к рассмотрению рукописей, не могут быть использованы в личных исследованиях без письменного согласия автора.

Информация или идеи, полученные в ходе рецензирования и связанные с возможными преимуществами, должны сохраняться конфиденциальными и не использоваться с целью получения личной выгоды.

\item[2.]

Все авторы обязаны раскрыть финансовые или другие явные или потенциальные конфликты интересов, которые могут быть восприняты как оказавшие влияние на результаты или выводы, представленные в работе.

\item[3.]

Рецензенты не должны участвовать в рассмотрении рукописей в случае наличия конфликтов интересов вследствие конкурентных, совместных и других взаимодействий и отношений с любым из авторов, компаниями или другими организациями, связанными с представленной работой.

\end{enumerate}



\noindent

{\bf 11. Положение о конфиденциальности}

\begin{enumerate}

\item[1.] Следующая информация об авторах, включая их имена, аффилиацию, контактные сведения, сведения о занимаемой должности и академической степени, ссылки на авторские профили, размещается в статье в разделе «Сведения об авторе».

\item[2.]

Паспортные данные и место жительство авторов, предоставленные в редакционную коллегию журнала, используются исключительно для заключения лицензионных договоров и пересылки корреспонденции и не передаются третьим лицам.

\item[3.]

Имена рецензентов не сообщаются авторам и третьим лицам.

\end{enumerate}





\noindent

{\bf 12. Положения о плагиате, дублировании и избыточности публикации}



\begin{enumerate}

\item[1.]  Каждая статья проверяется на наличие неправомерных заимствований из других работ (плагиат), дублирование и избыточность публикации (включая перевод публикации с других языков).

\item[2.]  

Для выявления плагиата редакционная коллегия использует следующие ресурсы: Антиплагиат (\url{http://samgtu.antiplagiat.ru/}) и Exactus Like (\url{http://like.exactus.ru/index.php/ru/}).

\item[3.]  

Статья, содержащая неправомерные заимствования из других работ (плагиат),  признаки дублирования и избыточности (в том числе признаки автоплагиата), является неприемлемой. Такая статья отклоняется без рецензирования.

\item[4.]  

Расширенные результаты, представленные на конференциях и семинарах, публикуемые в журнале по согласованию с редакционной коллегией и имеющие соответствующие ссылки, не считаются дублированной и избыточной публикацией.

\item[5.]  

Авторам разрешается подавать в журнал статьи, опубликованные на серверах препринтов, кроме опубликованных ранее в других изданиях. Такая статья не считается дублированной и избыточной.

\item[6.]  

Если статья опубликована, но позднее было обнаружено, что она содержит неправомерные заимствования из других работ (плагиат) или является дублированной, то редактор обращается к \href{http://publicationethics.org/resources/flowcharts}{блок-схемам} Комитета по издательской этике (\href{http://publicationethics.org/}{COPE}) для удаления такой статьи.

\end{enumerate}



\noindent

{\bf 13. Повторное использование текста}



\begin{enumerate}

\item[1.] 

Доля повторно использованного текста (текста из других авторских произведений) в представляемой в журнал работе не должна превышать 15\% от общего объёма статьи (не считая списка литературы).

\item[2.]

Повторное использование текста допустимо лишь для включения необходимых определений, формулировки основных принципов и т.п., чтобы читатель мог ознакомиться со статьей без привлечения дополнительных источников. В этом случае повторное использование текста не считается автоплагиатом.

\item[3.]

Если авторы считают необходимым ознакомить читателя с материалами, близкими к тематике статьи или оставшимися за рамками статьи, то они могут приводить достаточно полный список литературы, в которой интересующийся читатель может найти эти материалы.

\end{enumerate}



%\end{document}

